\scanpage{311}
\mspage{261}

Supposons $p_1\in W$. Alors dans la factorisation de $p_1$ en $X_1\mathrel{-}\overline X_{\unsure{1|2}}\mathrel{\unsure{-|\to}} S$, on peut prendre $\overline X_1=S$, donc le carré est $W$-cartésien ssi
\[
(X_3\longrightarrow\overline{\deleted{\unsure{S|X}}X}_1\times_S X_2=X_2)\ \in W
\]
i.e. $f_1\in W$, qed.

Bien sûr, on ne peut en général affirmer les implications inverses ; $f_i\in W$ n'implique nullement $p_i\in W$.

Tous ces résultats, et la notion de carré $W$-cartésien, se dualisent (moyennant l'axiome \inserted{idoine} de factorisation à droite, cf page 240) en une théorie des sommes amalgamées $W$-homotopiques, et des carrés $W$-cocartésiens. Cette fois on travaille sur la catégorie
\diagram{\(\displaystyle
\unsure{\Phi|\Psi}=\left(\begin{tikzcd}[ampersand replacement=\&,column sep=tiny,row sep=tiny]
\&1\arrow[dl]\\
0\&\\
\&2\arrow[ul]
\end{tikzcd}\right)
\overset{j}{\lhook\joinrel\longrightarrow}
\Phi^{\sim}=\left(\begin{tikzcd}[ampersand replacement=\&,column sep=tiny,row sep=tiny]
\&1\arrow[dl]\&\\
0\&\&3\arrow[ul]\arrow[dl]\\
\&2\arrow[ul,dash]\&
\end{tikzcd}\right),\)}
$j$ étant une immersion fermée. \unsure{On a|Ceci|Cela} un foncteur restriction
\[
j^*:\mathrm{Ho}_W(\Phi^{\sim})\mathrel{\unsure{-|\to}}\mathrm{Ho}_W(\Phi),
\]
dont on construit un adjoint à \unsure{gauche|droite}
\[
j_{\unsure{!|*}}:\mathrm{Ho}_W(\Phi)\mathrel{\unsure{-|\to}}\mathrm{Ho}_W(\Phi^{\sim}).
\]
on trouve \deleted{\unsure{une}} notion \unsure{naturelle} de
% archive sheet number: 311
